\documentclass{article}
\usepackage{iclr2026_conference,times}
\input{math_commands.tex}
\usepackage{booktabs}
\usepackage{graphicx}
\usepackage{hyperref}
\usepackage{url}

\title{Where Bits Matter in DINO-WM Planning: Consolidated Interim Results on Apple M3}
\author{}

\begin{document}
\maketitle
\lhead{Under review as a workshop paper at ICLR 2026 (ES-Reasoning)}

\begin{abstract}
This report consolidates all currently available experiment outputs for module-aware quantization in DINO-WM planning on Wall.
The active strict run provides 20 completed cells over two budgets (bA,bB), two seeds, and five variants (FP16, uniform INT8, mixed INT8, uniform INT4, mixed INT4), each evaluated on three paired episodes.
A prior pilot FP16 anchor with 10 episodes is also available.
Current evidence supports three points: (1) 8-bit variants remain close to FP16, (2) 4-bit performance is substantially lower overall, and (3) mixed INT4 does not yet show a uniform advantage across both budgets in the partial results.
The where-bits localization block and paired-delta post-processing are still pending.
\end{abstract}

\section{Introduction}
Our target claim is that in low-bit world-model planning, \emph{where} precision is allocated can matter as much as \emph{how many} bits are used.
This document is an interim consolidation, not a final claim paper.
It reports only completed artifacts and avoids extrapolation beyond observed outputs.

\section{Setup}
\paragraph{Task and model.}
All consolidated results use DINO-WM on Wall with paired-goal evaluation.

\paragraph{Hardware and runtime.}
Runs were executed on Apple Silicon M3 with MPS backend.

\paragraph{Strict run block (completed core).}
The strict block includes 20 completed cells:
\begin{itemize}
  \item Variants: \textit{fp16}, \textit{uniform\_int8}, \textit{mixed\_int8}, \textit{uniform\_int4}, \textit{mixed\_int4}
  \item Budgets: \textit{bA}, \textit{bB}
  \item Seeds: 1, 2
  \item Episodes per cell: 3
\end{itemize}
Total evaluated episodes in this block: 60.

\paragraph{Additional pilot artifact.}
One prior FP16 pilot metric exists (10 episodes, bA/seed1) from run \textit{m3\_main2exp\_v2}; we use it as context only because episode count differs.

\section{Consolidated Results}
\subsection{Core variant summary (20-cell strict block)}
\begin{figure}[t]
\centering
\includegraphics[width=0.98\linewidth]{../figures_m3_main2_expansion/ m3_main2exp_v3/interim_metrics/interim_frontier_size_vs_success.pdf}
\caption{Interim accuracy-memory frontier using completed strict-run metrics.}
\label{fig:interim_frontier}
\end{figure}

\begin{table}[t]
\centering
\small
\begin{tabular}{lcccc}
\toprule
Variant & Mean Success & Success Std & Mean Time (s) & Mean Size (MB) \\
\midrule
FP16 & 0.8333 & 0.1667 & 278.45 & 204.99 \\
Uniform INT8 & 0.7500 & 0.2764 & 301.33 & 87.72 \\
Mixed INT8 & 0.8333 & 0.1667 & 295.90 & 148.31 \\
Uniform INT4 & 0.1667 & 0.1667 & 294.38 & 68.12 \\
Mixed INT4 & 0.1667 & 0.2887 & 306.91 & 138.84 \\
\bottomrule
\end{tabular}
\caption{Aggregated means over completed strict-run cells (bA,bB; seeds 1,2; 3 episodes per cell).}
\label{tab:core_summary}
\end{table}

\subsection{Budget-conditioned mixed-vs-uniform INT4}
\begin{figure}[t]
\centering
\includegraphics[width=0.98\linewidth]{../figures_m3_main2_expansion/ m3_main2exp_v3/interim_metrics/interim_budget_variant_bars.pdf}
\caption{Budget-wise variant comparison with seed-level points (completed strict-run cells).}
\label{fig:interim_budget_bars}
\end{figure}

\begin{table}[t]
\centering
\small
\begin{tabular}{lccc}
\toprule
Budget & Uniform INT4 & Mixed INT4 & Delta (Mixed - Uniform) \\
\midrule
bA & 0.1667 & 0.3333 & +0.1667 \\
bB & 0.1667 & 0.0000 & -0.1667 \\
\bottomrule
\end{tabular}
\caption{Budget-wise means over seeds 1 and 2 in the completed strict block.}
\label{tab:budget_delta}
\end{table}

\begin{figure}[t]
\centering
\includegraphics[width=0.88\linewidth]{../figures_m3_main2_expansion/ m3_main2exp_v3/interim_metrics/interim_mixed_vs_uniform_int4_delta.pdf}
\caption{Direct per-budget delta for mixed INT4 vs uniform INT4 (positive means mixed is better).}
\label{fig:interim_delta}
\end{figure}

\paragraph{Pilot context (separate protocol).}
The prior FP16 pilot point (10 episodes) reports success 0.50 at bA/seed1, compared to 1.00 for the strict-run FP16 bA/seed1 cell (3 episodes).
This variance reinforces that the current sample size is still limited.

\section{Interpretation}
Current data already supports these bounded conclusions:
\begin{itemize}
  \item 8-bit settings remain strong in this setup and close to FP16.
  \item 4-bit settings are notably weaker than 8-bit/FP16 anchors.
  \item Mixed INT4 is not uniformly better than uniform INT4 in the partial strict results; its direction depends on budget.
\end{itemize}
Therefore, the final claim ``mixed INT4 $>$ uniform INT4'' remains unresolved until remaining blocks complete.

\section{Pending Work Before Final Claims}
\begin{itemize}
  \item Where-bits localization block: \textit{encfp16\_50} vs \textit{encheadfp16\_50}
  \item Aggregation and episode-level paired delta outputs
  \item Confidence intervals/bootstrap for mixed-vs-uniform comparisons
\end{itemize}

\section{Limitations}
This interim analysis is bounded by small episode counts (3 per strict cell), one environment (Wall), one checkpoint family, and incomplete post-processing.
Timing reflects this software path and should not be interpreted as optimized kernel-level benchmarking.

\section{Conclusion}
The consolidated interim evidence confirms that low-bit behavior is regime-dependent and that 4-bit performance is fragile.
However, with current partial outputs, the mixed-vs-uniform INT4 advantage is budget-sensitive rather than universally positive.
Final conclusions require completion of the pending where-bits and paired-delta stages.

\end{document}
